\chapter{2007年福建卷（理）}

\section{单选题}

\begin{problem}
复数 \(\frac{1}{(1 + \i)^2}\) 等于（\quad）

\choices
    {\( \frac{1}{2} \)}
    {\( -\frac{1}{2} \)}
    {\(\frac{1}{2}\i\)}
    {\(-\frac{1}{2}\i\)}
\end{problem}

\begin{answer}
D
\end{answer}

\begin{solution}
由
\[
(1+\i)^2=1+2\i+\i^2=2\i
\]
得
\[
\frac{1}{(1+\i)^2}=\frac{1}{2\i}=-\frac{1}{2}\i
\]
故选 D.
\end{solution}

\begin{problem}
数列 \(\{a_{n}\}\) 的前 \(n\) 项和为 \(S_{n}\)，若 \(a_{n} = \frac{1}{n(n + 1)}\)，则 \(S_{5}\) 等于（\quad）

\choices
    {\(1\)}
    {\( \frac{5}{6} \)}
    {\( \frac{1}{6} \)}
    {\(\frac{1}{30}\)}
\end{problem}

\begin{answer}
B
\end{answer}

\begin{solution}
因为
\[
\frac{1}{n(n+1)}=\frac{1}{n}-\frac{1}{n+1}
\]
所以
\[
S_5=\sum_{n=1}^{5}\left(\frac{1}{n}-\frac{1}{n+1}\right)=1-\frac{1}{6}=\frac{5}{6}
\]
故选 B.
\end{solution}

\begin{problem}
已知集合 \(A = \{x \mid x < a\}\)，\(B = \{x \mid 1 < x < 2\}\)，且 \(A \cup (\complement_{\R} B) = \R\)，则实数 \(a\) 的取值范围是（\quad）

\choices
    {\(a \leqslant 1\)}
    {\(a < 1\)}
    {\(a \geqslant 2\)}
    {\(a > 2\)}
\end{problem}

\begin{answer}
C
\end{answer}

\begin{solution}
由 \(A\cup(\complement_{\R}B)=\R\)，对任意 \(x\in B\)，因为 \(x\notin\complement_{\R}B\)，必有 \(x\in A\)，故 \(B\subset A\). 反过来，若 \(B\subset A\)，则 \(A\cup(\complement_{\R}B)=\R\). 因此题设等价于 \((1,2)\subset(-\infty,a)\). 这要求且只要求 \(a\geqslant2\)：当 \(a\geqslant2\) 时任意 \(x\in(1,2)\) 都满足 \(x<a\)；当 \(a<2\) 时可取充分接近 \(2\) 且大于 \(a\) 的 \(x\in(1,2)\)，此时 \(x\notin A\). 故选 C.
\end{solution}

\begin{problem}
对于向量 \(\bs{a}\)，\(\bs{b}\)，\(\bs{c}\) 和实数 \(\lambda\)，下列命题中真命题的是（\quad）

\choices
    {若 \(\bs{a} \cdot \bs{b} = 0\)，则 \(\bs{a} = 0\) 或 \(\bs{b} = 0\)}
    {若 \(\lambda\bs{a} = 0\)，则 \(\lambda = 0\) 或 \(\bs{a} = 0\)}
    {若 \(\bs{a}^2 = \bs{b}^2\)，则 \(\bs{a} = \bs{b}\) 或 \(\bs{a} = -\bs{b}\)}
    {若 \(\bs{a} \cdot \bs{b} = \bs{a} \cdot \bs{c}\)，则 \(\bs{b} = \bs{c}\)}
\end{problem}

\begin{answer}
B
\end{answer}

\begin{solution}
命题 A 不正确，例如取 \(\bs{a}=(1,0)\)，\(\bs{b}=(0,1)\)，则 \(\bs{a}\cdot\bs{b}=0\)，但 \(\bs{a}\)，\(\bs{b}\) 均不是零向量. 命题 B 正确，因为若 \(\lambda\bs{a}=0\)，当 \(\lambda\neq0\) 时两边同除以 \(\lambda\) 得 \(\bs{a}=0\)，所以必有 \(\lambda=0\) 或 \(\bs{a}=0\). 命题 C 不正确，例如 \(\bs{a}=(1,0)\)，\(\bs{b}=(0,1)\) 的长度相等，但 \(\bs{a}\neq\bs{b}\) 且 \(\bs{a}\neq-\bs{b}\). 命题 D 不正确，例如取 \(\bs{a}=(1,0)\)，\(\bs{b}=(0,0)\)，\(\bs{c}=(0,1)\)，则 \(\bs{a}\cdot\bs{b}=\bs{a}\cdot\bs{c}=0\)，但 \(\bs{b}\neq\bs{c}\). 故选 B.
\end{solution}

\begin{problem}
已知函数 \( f(x) = \sin \left( \omega x + \frac{\pi}{3} \right) \)（\(\omega > 0\)）的最小正周期为 \( \pi \)，则该函数的图象（\quad）

\choices
    {关于点 \(\left(\frac{\pi}{3}, 0\right)\) 对称}
    {关于直线 \(x = \frac{\pi}{4}\) 对称}
    {关于点 \(\left(\frac{\pi}{4}, 0\right)\) 对称}
    {关于直线 \(x = \frac{\pi}{3}\) 对称}
\end{problem}

\begin{answer}
A
\end{answer}

\begin{solution}
函数的最小正周期为 \(\dfrac{2\pi}{\omega}=\pi\)，所以 \(\omega=2\)，即 \(f(x)=\sin\left(2x+\dfrac{\pi}{3}\right)\). 其图象与 \(x\) 轴的交点满足
\(2x+\frac{\pi}{3}=k\pi\)，\(x=\frac{k\pi}{2}-\frac{\pi}{6}\)，\((k\in\Z)\).
取 \(k=1\)，得交点 \(\left(\dfrac{\pi}{3},0\right)\). 正弦曲线关于这些交点中心对称，因此选项 A 正确. 其对称轴满足 \(2x+\dfrac{\pi}{3}=\dfrac{\pi}{2}+k\pi\)，即 \(x=\dfrac{\pi}{12}+\dfrac{k\pi}{2}\)，既不为 \(\dfrac{\pi}{4}\)，也不为 \(\dfrac{\pi}{3}\)；而 \(f\left(\dfrac{\pi}{4}\right)=\dfrac12\neq0\)，选项 C 也不正确. 故选 A.
\end{solution}

\begin{problem}
以双曲线 \(\frac{x^2}{9} - \frac{y^2}{16} = 1\) 的右焦点为圆心，且与其渐近线相切的圆的方程是（\quad）

\choices
    {\(x^{2} + y^{2} - 10x + 9 = 0\)}
    {\(x^{2} + y^{2} - 10x + 16 = 0\)}
    {\(x^{2} + y^{2} + 10x + 16 = 0\)}
    {\(x^{2} + y^{2} + 10x + 9 = 0\)}
\end{problem}

\begin{answer}
A
\end{answer}

\begin{solution}
双曲线的半实轴长和半虚轴长分别为 \(3\) 和 \(4\)，所以焦半距为
\[
c=\sqrt{3^2+4^2}=5
\]
右焦点为 \(F(5,0)\)，其渐近线为 \(y=\pm\dfrac{4}{3}x\). 取渐近线 \(4x-3y=0\)，焦点到该直线的距离为
\[
r=\frac{|4\times5-3\times0|}{\sqrt{4^2+(-3)^2}}=4
\]
因此圆的圆心为 \((5,0)\)，半径为 \(4\)，方程为
\[
(x-5)^2+y^2=4^2
\]
即 \(x^2+y^2-10x+9=0\). 故选 A.
\end{solution}

\begin{problem}
已知 \( f(x) \) 为 \( \R \) 上的减函数，则满足 \( f\left(\left|\frac{1}{x}\right|\right) < f(1) \) 的实数 \( x \) 的取值范围是（\quad）

\choices
    {\((-1, 1)\)}
    {\((0,1)\)}
    {\((-1,0)\cup (0,1)\)}
    {\((-\infty, -1) \cup (1, +\infty)\)}
\end{problem}

\begin{answer}
C
\end{answer}

\begin{solution}
函数 \(f\) 在 \(\R\) 上为严格减函数，所以
\[
f\left(\left|\frac{1}{x}\right|\right)<f(1)
\Longleftrightarrow \left|\frac{1}{x}\right|>1
\]
这里必须有 \(x\neq0\)，且 \(\left|\dfrac{1}{x}\right|>1\) 等价于 \(0<|x|<1\). 因而 \(x\) 的取值范围为 \((-1,0)\cup(0,1)\)，故选 C.
\end{solution}

\begin{problem}
已知 \(m\)，\(n\) 为两条不同的直线，\(\alpha\)，\(\beta\) 为两个不同的平面，则下列命题中正确的是（\quad）

\choices
    {\(m \subset \alpha\)，\(n \subset \alpha\)，\(m \parallel \beta\)，\(n \parallel \beta \Rightarrow \alpha \parallel \beta\)}
    {\(\alpha \parallel \beta\)，\(m \subset \alpha\)，\(n \subset \beta \Rightarrow m \parallel n\)}
    {\(m \perp \alpha\)，\(m \perp n \Rightarrow n \parallel \alpha\)}
    {\(n \parallel m\)，\(n \subset \alpha \Rightarrow m \perp \alpha\)}
\end{problem}

\begin{answer}
题面严格定义下无正确选项；公开答案存在 C/D 分歧
\end{answer}

\begin{solution}
逐项检验. 命题 A 不成立：取相交平面 \(\alpha\colon y=0\)，\(\beta\colon z=0\)，在 \(\alpha\) 内取两条沿 \(x\) 轴方向、分别位于 \(z=1\) 与 \(z=2\) 的直线，则它们都与 \(\beta\) 无公共点而平行于 \(\beta\)，但 \(\alpha\) 与 \(\beta\) 相交. 命题 B 不成立：取 \(\alpha\colon z=0\)，\(\beta\colon z=1\)，在 \(\alpha\) 内取沿 \(x\) 轴方向的直线 \(m\)，在 \(\beta\) 内取沿 \(y\) 轴方向的直线 \(n\)，虽有 \(\alpha\parallel\beta\)，但 \(m\not\parallel n\).

命题 C 也不成立. 取 \(\alpha\colon z=0\)，取 \(m\) 为 \(z\) 轴，取 \(n\) 为 \(x\) 轴，则 \(m\perp\alpha\) 且 \(m\perp n\)，但 \(n\subset\alpha\)，按直线与平面平行的通常定义，\(n\) 不平行于 \(\alpha\). 命题 D 不成立：仍取 \(\alpha\colon z=0\)，令 \(n\colon y=z=0\)，令 \(m\colon y=1\)，\(\ z=1\)，则 \(n\parallel m\)、\(n\subset\alpha\)，但 \(m\) 的方向平行于 \(\alpha\)，故 \(m\not\perp\alpha\). 因此按原 PDF 的题面和通常定义，四个命题均不成立，题面没有正确选项；现有公开答案分别出现 C、D，不能据此伪造唯一数学结论.
\end{solution}

\begin{problem}
把 \(1 + (1 + x) + (1 + x)^2 + \cdots + (1 + x)^n\) 展开成关于 \(x\) 的多项式，其各项系数和为 \(a_n\)，则 \(\lim\limits_{n \to \infty} \frac{2a_n - 1}{a_n + 1}\) 等于（\quad）

\choices
    {\( \frac{1}{4} \)}
    {\( \frac{1}{2} \)}
    {\(1\)}
    {\(2\)}
\end{problem}

\begin{answer}
D
\end{answer}

\begin{solution}
多项式各项系数之和等于令 \(x=1\) 时的值，因此
\[
a_n=1+2+2^2+\cdots+2^n=2^{n+1}-1
\]
于是
\[
\lim_{n\to\infty}\frac{2a_n-1}{a_n+1}
 =\lim_{n\to\infty}\frac{2^{n+2}-3}{2^{n+1}}=2
\]
故选 D.
\end{solution}

\begin{problem}
顶点在同一球面上的正四棱柱 \(ABCD - A'B'C'D'\) 中，\(AB = 1\)，\(AA' = \sqrt{2}\)，则 \(A\)，\(C\) 两点间的球面距离为（\quad）

\choices
    {\( \frac{\pi}{4} \)}
    {\( \frac{\pi}{2} \)}
    {\( \frac{\sqrt{2}}{4}\pi \)}
    {\( \frac{\sqrt{2}}{2}\pi \)}
\end{problem}

\begin{answer}
B
\end{answer}

\begin{solution}
设该正四棱柱外接球的球心为 \(O\). 正四棱柱的空间对角线为
\[
AC'=\sqrt{AB^2+BC^2+AA'^2}=\sqrt{1+1+2}=2
\]
故外接球半径为 \(R=1\). 设球心角 \(\angle AOC=\theta\)，则弦长 \(AC=\sqrt2\) 满足
\[
\sqrt2=2R\sin\frac{\theta}{2}=2\sin\frac{\theta}{2}
\]
由于 \(0\leqslant\theta\leqslant\pi\)，得 \(\theta=\dfrac\pi2\). 球面距离等于半径乘以对应的圆心角，所以所求距离为 \(R\theta=\dfrac\pi2\)，故选 B.
\end{solution}

\begin{problem}
已知对任意实数 \(x\)，有 \(f(-x) = -f(x)\)，\(g(-x) = g(x)\)，且 \(x > 0\) 时，\(f'(x) > 0\)，\(g'(x) > 0\)，则 \(x < 0\) 时，（\quad）

\choices
    {\(f'(x) > 0\)，\(g'(x) > 0\)}
    {\(f^{\prime}(x) > 0\)，\(g^{\prime}(x) < 0\)}
    {\(f^{\prime}(x) < 0\)，\(g^{\prime}(x) > 0\)}
    {\(f^{\prime}(x) < 0\)，\(g^{\prime}(x) < 0\)}
\end{problem}

\begin{answer}
B
\end{answer}

\begin{solution}
由 \(f(-x)=-f(x)\)，对等式两边求导得 \(f'(-x)=f'(x)\)，所以 \(f'\) 为偶函数. 由 \(g(-x)=g(x)\)，求导得 \(g'(-x)=-g'(x)\)，所以 \(g'\) 为奇函数. 当 \(x<0\) 时，\(-x>0\)，故
故 \(f'(x)=f'(-x)>0\)，\(g'(x)=-g'(-x)<0\).
因此选 B.
\end{solution}

\begin{problem}
将 \(9\) 个数 \(a_{ij}\)（\(i = 1\)，\(2\)，\(3\)，\(j = 1\)，\(2\)，\(3\)）排成三行三列的方阵，从中任取三个数，则至少有两个数位于同行或同列的概率是（\quad）

\choices
    {\(\frac{3}{7}\)}
    {\(\frac{4}{7}\)}
    {\(\frac{1}{14}\)}
    {\(\frac{13}{14}\)}
\end{problem}

\begin{answer}
D
\end{answer}

\begin{solution}
从 \(9\) 个数中任取 \(3\) 个，共有
\[
\mathrm C_9^3=84
\]
种取法. 若任取的三个数没有两个位于同行或同列，则三个数必须分别位于三行，且也分别位于三列. 这相当于把三列与三行一一对应，共有 \(3!=6\) 种取法. 所以至少有两个数位于同行或同列的取法数为 \(84-6=78\)，其概率为
\[
\frac{78}{84}=\frac{13}{14}
\]
故选 D.
\end{solution}

\section{填空题}

\begin{problem}
已知实数 \(x\)，\(y\) 满足 \(\begin{cases}
x + y \geqslant 2,\\
x - y \leqslant 2,\\
0 \leqslant y \leqslant 3,
\end{cases}\) 则 \(z = 2x - y\) 的取值范围是\fillinblank{}.
\end{problem}

\begin{answer}
\(\left[-5, 7\right]\)
\end{answer}

\begin{solution}
由题设得
\(x\geqslant2-y\)，\(x\leqslant2+y\)，\(0\leqslant y\leqslant3\)
对固定的 \(y\)，因为 \(z=2x-y\) 随 \(x\) 增大而增大，所以
\[
2(2-y)-y\leqslant z\leqslant2(2+y)-y
\]
即
\[
4-3y\leqslant z\leqslant4+y
\]
当 \(y=3\) 且 \(x=2-y=-1\) 时，取得最小值 \(z=-5\)；当 \(y=3\) 且 \(x=2+y=5\) 时，取得最大值 \(z=7\). 可行域为凸集，线性函数 \(z=2x-y\) 的值域为区间，因此 \(z\) 的取值范围是 \([-5,7]\).
\end{solution}

\begin{problem}
已知正方形 \(ABCD\)，则以 \(A\)，\(B\) 为焦点，且过 \(C\)，\(D\) 两点的椭圆的离心率为\fillinblank{}.
\end{problem}

\begin{answer}
\(\sqrt{2} - 1\)
\end{answer}

\begin{solution}
设正方形边长为 \(s\)，椭圆的半焦距为 \(c\)，半长轴为 \(a\). 由于焦点为相邻顶点 \(A\)，\(B\)，有
\(2c=AB=s\)，\(c=\frac{s}{2}\)
点 \(C\) 在椭圆上，且 \(CA=s\sqrt2\)，\(CB=s\)，所以椭圆定义给出
\[
2a=CA+CB=s\sqrt2+s=s(\sqrt2+1)
\]
因此离心率为
\[
e=\frac{c}{a}=\frac{1}{\sqrt2+1}=\sqrt2-1
\]
又 \(DA=BC=s\)，\(DB=AC=s\sqrt2\)，所以 \(DA+DB=s(\sqrt2+1)=2a\)，故点 \(D\) 也在该椭圆上. 因此所求离心率为 \(\sqrt2-1\).
\end{solution}

\begin{problem}
两封信随机投入 \(A\)，\(B\)，\(C\) 三个空邮箱，则 \(A\) 邮箱的信件数 \(\xi\) 的数学期望 \(E(\xi) =\fillinblank{}\).
\end{problem}

\begin{answer}
\(\frac{2}{3}\)
\end{answer}

\begin{solution}
设两封信分别为第一个和第二个投递结果，令 \(I_1\)，\(I_2\) 分别表示该封信投入邮箱 \(A\) 的示性变量，则 \(\xi=I_1+I_2\). 每封信投入 \(A\) 的概率均为 \(\dfrac13\)，所以
\[
E(I_1)=E(I_2)=\frac13
\]
由数学期望的线性性质，
\[
E(\xi)=E(I_1)+E(I_2)=\frac13+\frac13=\frac23
\]
\end{solution}

\begin{problem}
中学数学中存在许多关系，比如“相等关系”，“平行关系”等等. 如果集合 \(A\) 中元素之间的一个关系“\(\sim\)”满足以下三个条件：

\begin{circlelist}
\item 自反性：对于任意 \(a \in A\)，都有 \(a \sim a\)；
\item 对称性：对于 \(a\)，\(b \in A\)，若 \(a \sim b\)，则有 \(b \sim a\)；
\item 传递性：对于 \(a\)，\(b\)，\(c \in A\)，若 \(a \sim b\)，\(b \sim c\)，则有 \(a \sim c\).
\end{circlelist}

则称“\(\sim\)”是集合 \(A\) 的一个等价关系. 例如：“数的相等”是等价关系，而“直线的平行”不是等价关系（自反性不成立）. 请你再列出三个等价关系：\fillinblank{}.
\end{problem}

\begin{answer}
图形的全等、图形的相似、非零向量的共线
\end{answer}

\begin{solution}
图形的全等关系具有自反性、对称性和传递性：任一图形与自身全等；若图形 \(X\) 与图形 \(Y\) 全等，则图形 \(Y\) 与图形 \(X\) 全等；若图形 \(X\) 与图形 \(Y\) 全等且图形 \(Y\) 与图形 \(Z\) 全等，则图形 \(X\) 与图形 \(Z\) 全等. 因而图形的全等是等价关系. 图形的相似关系也满足三个条件：任一图形与自身相似；若图形 \(X\) 与图形 \(Y\) 相似，则图形 \(Y\) 与图形 \(X\) 相似；若图形 \(X\) 与图形 \(Y\) 相似且图形 \(Y\) 与图形 \(Z\) 相似，则图形 \(X\) 与图形 \(Z\) 相似. 因而图形的相似也是等价关系. 对非零向量，若 \(\bs{a}\sim\bs{b}\) 表示两向量共线，则 \(\bs{a}=\lambda\bs{b}\)，其中 \(\lambda\neq0\)，所以自反性成立；又有 \(\bs{b}=\lambda^{-1}\bs{a}\)，所以对称性成立；若 \(\bs{a}=\lambda\bs{b}\)，\(\bs{b}=\mu\bs{c}\)，则 \(\bs{a}=\lambda\mu\bs{c}\)，且 \(\lambda\mu\neq0\)，故传递性也成立. 所以图形的全等、图形的相似、非零向量的共线均为等价关系.
\end{solution}

\section{解答题}

\begin{problem}
在 \(\triangle ABC\) 中，\(\tan A = \frac{1}{4}\)，\(\tan B = \frac{3}{5}\).

\begin{enumerate}
\item 求角 \(C\) 的大小；
\item 若 \(AB\) 边的长为 \(\sqrt{17}\)，求 \(BC\) 边的长.
\end{enumerate}
\end{problem}

\begin{solution}
\begin{enumerate}
\item 由于 \(A\)，\(B\) 是三角形的内角，且 \(\tan A\)，\(\tan B\) 均为正数，所以 \(A\)，\(B\) 均为锐角. 因而
\[
\tan(A+B)=\frac{\tan A+\tan B}{1-\tan A\tan B}
=\frac{\frac{1}{4}+\frac{3}{5}}{1-\frac{1}{4}\cdot\frac{3}{5}}=1
\]
又 \(0<A+B<\pi\)，且 \(\tan(A+B)=1>0\)，所以 \(A+B\) 不可能在 \((\frac{\pi}{2},\pi)\) 内，只能为锐角. 因而 \(A+B=\frac{\pi}{4}\)，从而
\[
C=\pi-A-B=\frac{3\pi}{4}
\]

\item 由 \(\tan A=\frac{1}{4}\) 及 \(A\) 为锐角，得
\(\sin A=\frac{1}{\sqrt{17}}\)，\(\sin C=\sin\frac{3\pi}{4}=\frac{\sqrt{2}}{2}\).
由正弦定理，
\[
\frac{BC}{\sin A}=\frac{AB}{\sin C}
\]
所以
\[
BC=\sqrt{17}\cdot\frac{\sin A}{\sin C}
=\sqrt{17}\cdot\frac{1/\sqrt{17}}{\sqrt{2}/2}=\sqrt{2}
\]
故 \(BC=\sqrt{2}\).
\end{enumerate}
\end{solution}

\begin{problem}
如图，正三棱柱 \(ABC - A_{1}B_{1}C_{1}\) 的所有棱长都为 \(2\)，\(D\) 为 \(CC_{1}\) 中点.

\begin{enumerate}
\item 求证：\(AB_{1} \perp\) 平面 \(A_{1}BD\)；
\item 求二面角 \(A - AD_{1} - B\) 的大小；
\item 求点 \(C\) 到平面 \(A_{1}BD\) 的距离.
\end{enumerate}

\bitmapfigure[max width=0.42\linewidth]{img/2007/fujian_science/q18_fig1.png}
\end{problem}

\begin{solution}
建立空间直角坐标系，取点 \(A(0,0,0)\)，\(B(2,0,0)\)，\(C(1,\sqrt{3},0)\)，\(A_{1}(0,0,2)\)，\(B_{1}(2,0,2)\)，\(C_{1}(1,\sqrt{3},2)\).
由 \(D\) 为 \(CC_{1}\) 的中点，得 \(D(1,\sqrt{3},1)\).

\begin{enumerate}
\item 由坐标可得
\(\overrightarrow{AB_{1}}=(2,0,2)\)，\(\overrightarrow{A_{1}B}=(2,0,-2)\)，\(\overrightarrow{A_{1}D}=(1,\sqrt{3},-1)\). 于是 \(\overrightarrow{AB_{1}}\cdot\overrightarrow{A_{1}B}=4-4=0\)，\(\overrightarrow{AB_{1}}\cdot\overrightarrow{A_{1}D}=2-2=0\).
直线 \(A_{1}B\)，\(A_{1}D\) 是平面 \(A_{1}BD\) 内相交的两条直线，所以
直线 \(AB_{1}\) 垂直于平面 \(A_{1}BD\).

\item 原题第（2）问印作二面角 \(A-AD_{1}-B\)，但题干和图中均未定义 \(D_{1}\). 结合图中已定义的点，以下按二面角 \(A-A_{1}D-B\) 计算，其棱为 \(A_{1}D\)，两侧面为 \(AA_{1}D\) 和 \(A_{1}BD\).

平面 \(AA_{1}D\) 和平面 \(A_{1}BD\) 的一个法向量分别可取 \(\bs{n}_{1}=\overrightarrow{AA_{1}}\times\overrightarrow{AD}=(-2\sqrt{3},2,0)\)，\(\bs{n}_{2}=\overrightarrow{A_{1}B}\times\overrightarrow{A_{1}D}=(2\sqrt{3},0,2\sqrt{3})\). 因为 \(|\bs{n}_{1}|=4\)，\(|\bs{n}_{2}|=2\sqrt{6}\)，且 \(\bs{n}_{1}\cdot\bs{n}_{2}=-12\)，
所以二面角的大小 \(\theta\) 满足
\[
\cos\theta=\frac{|\bs{n}_{1}\cdot\bs{n}_{2}|}{|\bs{n}_{1}|\,|\bs{n}_{2}|}
=\frac{12}{4\cdot2\sqrt{6}}=\frac{\sqrt{6}}{4}
\]
故 \(\theta=\arccos\frac{\sqrt{6}}{4}\).

\item 由
\[
\overrightarrow{A_{1}B}\times\overrightarrow{A_{1}D}=(2\sqrt{3},0,2\sqrt{3})
\]
可取平面 \(A_{1}BD\) 的法向量为 \((1,0,1)\). 该平面过点 \(A_{1}(0,0,2)\)，故其方程为
\[
x+z-2=0
\]
于是点 \(C(1,\sqrt{3},0)\) 到平面 \(A_{1}BD\) 的距离为
\[
d=\frac{|1+0-2|}{\sqrt{1^{2}+0^{2}+1^{2}}}=\frac{\sqrt{2}}{2}
\]
\end{enumerate}
\end{solution}

\begin{problem}
某分公司经销某种品牌产品，每件产品的成本为 \(3\text{ 元}\)，并且每件产品需向总公司交 \(a\text{ 元}\)（\(3 \leqslant a \leqslant 5\)）的管理费，预计当每件产品的售价为 \(x\text{ 元}\)（\(9 \leqslant x \leqslant 11\)）时，一年的销售量为 \((12 - x)^{2}\text{ 万件}\).

\begin{enumerate}
\item 求分公司一年的利润 \(L\)（万元）与每件产品的售价 \(x\) 的函数关系式；
\item 当每件产品的售价为多少元时，分公司一年的利润 \(L\) 最大？并求出 \(L\) 的最大值 \(Q(a)\).
\end{enumerate}
\end{problem}

\begin{solution}
\begin{enumerate}
\item 每件产品的利润为售价减去成本和管理费，即 \(x-3-a\) 元. 销售量为 \((12-x)^{2}\) 万件，所以年利润为 \(L=(x-a-3)(12-x)^{2}\)，其中 \(9\leqslant x\leqslant 11\).

\item 对 \(L\) 求导，得
\[
L'=(12-x)^{2}-2(x-a-3)(12-x)=(12-x)(18+2a-3x)
\]
在 \(9\leqslant x\leqslant 11\) 时，\(12-x>0\)，所以极值只由 \(18+2a-3x\) 的符号决定. 令
\[
x_{0}=\frac{18+2a}{3}=6+\frac{2a}{3}
\]
当 \(3\leqslant a\leqslant\frac{9}{2}\) 时，\(x_{0}\leqslant 9\)，从而 \(L\) 在 \([9,11]\) 上递减，最大值在 \(x=9\) 处取得，并且
\[
Q(a)=L(9)=(6-a)\cdot 3^{2}=54-9a
\]

当 \(\frac{9}{2}<a\leqslant 5\) 时，\(9<x_{0}<11\)，\(L\) 在 \([9,x_{0}]\) 上递增，在 \([x_{0},11]\) 上递减，最大值在 \(x=x_{0}\) 处取得. 此时 \(x_{0}-a-3=3-\frac{a}{3}=\frac{9-a}{3}\)，\(12-x_{0}=6-\frac{2a}{3}=\frac{2(9-a)}{3}\)，
所以
\[
Q(a)=L(x_{0})=\frac{9-a}{3}\left(\frac{2(9-a)}{3}\right)^{2}
=\frac{4(9-a)^{3}}{27}
\]
综上，售价为
\[
x=\begin{cases}
9,&3\leqslant a\leqslant\frac{9}{2},\\
6+\frac{2a}{3},&\frac{9}{2}<a\leqslant5,
\end{cases}
\]
最大利润为
\[
Q(a)=\begin{cases}
54-9a,&3\leqslant a\leqslant\frac{9}{2},\\
\frac{4(9-a)^{3}}{27},&\frac{9}{2}<a\leqslant5
\end{cases}
\]
\end{enumerate}
\end{solution}

\begin{problem}
如图，已知点 \(F(1,0)\)，直线 \(l\)：\(x = -1\)，\(P\) 为平面上的动点. 过 \(P\) 作直线 \(l\) 的垂线，垂足为点 \(Q\)，且 \(\overrightarrow{QP} \cdot \overrightarrow{QF} = \overrightarrow{FP} \cdot \overrightarrow{FQ}\).

\begin{enumerate}
\item 求动点 \(P\) 的轨迹 \(C\) 的方程；
\item 过点 \(F\) 的直线交轨迹 \(C\) 于 \(A\)，\(B\) 两点，交直线 \(l\) 于点 \(M\)，已知 \(\overrightarrow{MA} = \lambda_1 \overrightarrow{AF}\)，\(\overrightarrow{MB} = \lambda_2 \overrightarrow{BF}\)，求 \(\lambda_1 + \lambda_2\) 的值.
\end{enumerate}

\bitmapfigure[max width=0.30\linewidth]{img/2007/fujian_science/q20_fig1.png}
\end{problem}

\begin{solution}
\begin{enumerate}
\item 设 \(P(x,y)\)，则垂足为 \(Q(-1,y)\). 因而 \(\overrightarrow{QP}=(x+1,0)\)，\(\overrightarrow{QF}=(2,-y)\)，\(\overrightarrow{FP}=(x-1,y)\)，\(\overrightarrow{FQ}=(-2,y)\).
代入题设向量等式，得
\[
2(x+1)=-2(x-1)+y^{2}
\]
化简得 \(y^{2}=4x\)，所以动点 \(P\) 的轨迹 \(C\) 的方程为
\[
y^{2}=4x
\]

\item 设过 \(F\) 的直线方向向量为 \((u,v)\)，用参数 \(t\) 表示直线上各点：
\[
X(t)=(1+tu,tv)
\]
由于该直线与 \(l\colon x=-1\) 相交，故 \(u\neq0\). 设交点 \(M=X(t_{0})\)，则由 \(1+t_{0}u=-1\)，得 \(t_{0}=-\frac{2}{u}\).
直线与抛物线交于 \(A=X(t_{1})\)，\(B=X(t_{2})\). 代入 \(y^{2}=4x\)，得
\[
v^{2}t^{2}-4ut-4=0
\]
因为有两个交点，所以 \(v\neq0\). 由韦达定理，\(t_{1}+t_{2}=\frac{4u}{v^{2}}\)，\(t_{1}t_{2}=-\frac{4}{v^{2}}\)，从而 \(\frac{t_{1}+t_{2}}{t_{1}t_{2}}=-u\).
又 \(\overrightarrow{MA}=(t_{1}-t_{0})(u,v)\)，\(\overrightarrow{AF}=-t_{1}(u,v)\).
由 \(\overrightarrow{MA}=\lambda_{1}\overrightarrow{AF}\)，得
\[
\lambda_{1}=\frac{t_{0}}{t_{1}}-1
\]
由 \(\overrightarrow{MB}=(t_{2}-t_{0})(u,v)\)，\(\overrightarrow{BF}=-t_{2}(u,v)\)，代入 \(\overrightarrow{MB}=\lambda_{2}\overrightarrow{BF}\)，得
\[
\lambda_{2}=\frac{t_{0}}{t_{2}}-1
\]
因此
\[
\lambda_{1}+\lambda_{2}
=t_{0}\left(\frac{1}{t_{1}}+\frac{1}{t_{2}}\right)-2
=t_{0}\frac{t_{1}+t_{2}}{t_{1}t_{2}}-2
=\left(-\frac{2}{u}\right)(-u)-2=0
\]
故 \(\lambda_{1}+\lambda_{2}=0\).
\end{enumerate}
\end{solution}

\begin{problem}
等差数列 \(\{a_{n}\}\) 的前 \(n\) 项和为 \(S_{n}\)，\(a_{1} = 1 + \sqrt{2}\)，\(S_{3} = 9 + 3\sqrt{2}\).

\begin{enumerate}
\item 求数列 \(\{a_{n}\}\) 的通项 \(a_{n}\) 与前 \(n\) 项和为 \(S_{n}\)；
\item 设 \(b_{n} = \frac{S_{n}}{n}\)（\(n \in \N^*\)），求证：数列 \(\{b_{n}\}\) 中任意不同的三项都不可能成为等比数列.
\end{enumerate}
\end{problem}

\begin{solution}
\begin{enumerate}
\item 设等差数列 \(\{a_{n}\}\) 的公差为 \(d\). 由等差数列的前 \(3\) 项和公式，
\[
S_{3}=3(a_{1}+d)
\]
所以
\[
3(1+\sqrt{2}+d)=9+3\sqrt{2}
\]
从而 \(d=2\).
\[
a_{n}=a_{1}+(n-1)d=2n-1+\sqrt{2}
\]
且
\[
S_{n}=\frac{n(a_{1}+a_{n})}{2}=n(n+\sqrt{2})
\]

\item 由上式，
\[
b_{n}=\frac{S_{n}}{n}=n+\sqrt{2}
\]
任取三个不同的正整数 \(n_{1}<n_{2}<n_{3}\)，令 \(p=n_{2}-n_{1}>0\)，\(q=n_{3}-n_{2}>0\).
因为 \(b_{n}\) 随 \(n\) 严格递增，若这三项能组成等比数列，则中间项必须满足
\[
b_{n_{2}}^{2}=b_{n_{1}}b_{n_{3}}
\]
将 \(b_{n}=n+\sqrt{2}\) 代入，记 \(t=n_{1}+\sqrt{2}\)，则该等式等价于
\[
(t+p)^{2}=t(t+p+q)
\]
即
\[
t(p-q)+p^{2}=0
\]
若 \(p=q\)，则上式变为 \(p^{2}=0\)，与 \(p>0\) 矛盾；若 \(p\neq q\)，则
\[
\sqrt{2}=-n_{1}-\frac{p^{2}}{p-q}
\]
为有理数，也矛盾. 因此任意不同的三项都不可能成为等比数列.
\end{enumerate}
\end{solution}

\begin{problem}
已知函数 \(f(x) = \e^{x} - kx\)，\(x \in \R\).

\begin{enumerate}
\item 若 \(k = \e\)，试确定函数 \(f(x)\) 的单调区间；
\item 若 \(k > 0\)，且对于任意 \(x \in \R\)，\(f(|x|) > 0\) 恒成立，试确定实数 \(k\) 的取值范围；
\item 设函数 \(F(x) = f(x) + f(-x)\)，求证：\(F(1)F(2)\cdots F(n) > (\e^{n+1} + 2)^{\frac{n}{2}}\)（\(n \in \N^*\)）.
\end{enumerate}
\end{problem}

\begin{solution}
\begin{enumerate}
\item 当 \(k=\e\) 时，
\[
f'(x)=\e^{x}-\e
\]
因为 \(\e^{x}<\e\) 当且仅当 \(x<1\)，所以 \(f'(x)<0\) 在 \((-\infty,1)\) 上成立，\(f'(x)>0\) 在 \((1,+\infty)\) 上成立. 因此 \(f(x)\) 在 \((-\infty,1]\) 上单调递减，在 \([1,+\infty)\) 上单调递增.

\item 令 \(t=|x|\)，则 \(t\geqslant0\)，题设等价于
\[
\e^{t}-kt>0\quad\text{对任意 }t\geqslant0
\]
当 \(t=0\) 时，\(\e^{t}-kt=1>0\). 当 \(t>0\) 时，不等式等价于
\[
k<\frac{\e^{t}}{t}
\]
令 \(h(t)=\frac{\e^{t}}{t}\)，则
\[
h'(t)=\frac{\e^{t}(t-1)}{t^{2}}
\]
所以 \(h(t)\) 在 \((0,1)\) 上递减，在 \((1,+\infty)\) 上递增，最小值为 \(h(1)=\e\). 因为题设是不等式 \(f(|x|)>0\)，在 \(t=1\) 处也必须严格成立，故
\[
0<k<\e
\]

\item 由函数定义，
\[
F(x)=\e^{x}-kx+\e^{-x}+kx=\e^{x}+\e^{-x}
\]
记 \(P_{n}=F(1)F(2)\cdots F(n)\). 对 \(1\leqslant j\leqslant\left\lfloor\frac{n}{2}\right\rfloor\)，有
\[
F(j)F(n+1-j)
=\e^{n+1}+\e^{-(n+1)}+\e^{n+1-2j}+\e^{2j-(n+1)}
\]
由于最后两项的乘积为 \(1\)，且 \(j\neq\frac{n+1}{2}\)，由 \(u+u^{-1}>2\)（\(u>0\)，\(u\neq1\)）得
\[
F(j)F(n+1-j)>\e^{n+1}+2
\]

若 \(n\) 为偶数，共有 \(n/2\) 对上述因子，所以
\[
P_{n}>\left(\e^{n+1}+2\right)^{\frac{n}{2}}
\]

若 \(n\) 为奇数，除上述 \((n-1)/2\) 对因子外，还有中间项 \(F\left(\frac{n+1}{2}\right)\). 且
\[
F\left(\frac{n+1}{2}\right)^{2}
=\e^{n+1}+2+\e^{-(n+1)}>\e^{n+1}+2
\]
从而
\[
F\left(\frac{n+1}{2}\right)>\left(\e^{n+1}+2\right)^{\frac12}
\]
因此
\[
P_{n}>\left(\e^{n+1}+2\right)^{\frac{n-1}{2}}
\left(\e^{n+1}+2\right)^{\frac12}
=\left(\e^{n+1}+2\right)^{\frac{n}{2}}
\]
两种情形均得所证不等式.
\end{enumerate}
\end{solution}
